\chapter{Risk, Operations, Product Control and Compliance}\label{ch:in:risk-operations-product-control-and-compliance}

The European Banking Authority's guidelines on pay say that ``The remuneration of independent \omterm{def:in:risk-operations-product-control-and-compliance:control}{control functions} should
be predominantly fixed, to reflect the nature of their responsibilities'', and that any variable pay should be set
``separately from the business units they control''. The rule shapes the career as much as the pay: a risk manager,
\omterm{def:in:risk-operations-product-control-and-compliance:controller}{product controller} or \omterm{def:in:risk-operations-product-control-and-compliance:compliance}{compliance officer} is paid for the quality of the control, not for the profit it allows. The
\omterm{def:in:risk-operations-product-control-and-compliance:control}{control functions} employ many of the industry's quantitative people, and for many of them they are a route into the
industry as often as a destination. This chapter describes the four functions, compares their pay with the front
office's in the filings and in the banks' reports, and follows the routes between them.

\begin{center}\footnotesize
\begin{tabular}{@{}p{2.2cm}p{2.5cm}p{2.5cm}p{2.5cm}p{2.5cm}@{}}
\toprule
\multicolumn{5}{@{}l}{\textbf{Role cards: the control and support functions}}\\
\midrule
 & risk manager & operations analyst & \omterm{def:in:risk-operations-product-control-and-compliance:controller}{product controller} & \omterm{def:in:risk-operations-product-control-and-compliance:compliance}{compliance officer}\\
\midrule
works on & limits, measures, escalation & settlement, confirmations, reconciliation & daily P\&L and valuation & rules, surveillance, advice\\
reports to & chief risk officer & chief operating officer & chief financial officer & head of compliance\\
codes in filings & 13-2054, 13-2099.01 & 43-4011, 13-2099 & 13-2011, 13-2051 & 13-1041, 23-1011\\
taught in & Book~6, ch.~21; Book~16, ch.~12 & Book~16, ch.~13 & Book~6, ch.~27; Book~16, ch.~15 & Book~16, ch.~16\\
\bottomrule
\end{tabular}
\end{center}

\begin{definition}[Control function]\label{def:in:risk-operations-product-control-and-compliance:control}
A \emph{control function}\index{control function} is an organisational unit, independent of the businesses it
oversees, whose job is to measure, check and limit what those businesses do: risk management, compliance and internal
audit in the regulators' list, and in trading firms also product control, which checks the profit and loss the front
office reports.
\end{definition}

The EBA's guidelines say independent \omterm{def:in:risk-operations-product-control-and-compliance:control}{control functions} ``typically comprise risk management, compliance and internal
audit functions''. Operations is not a \omterm{def:in:risk-operations-product-control-and-compliance:control}{control function} in that sense, but it is independent of the front office in the
same way (the front, middle and back offices of Book~16, chapter~13) and it employs many of the same people.

\section{Risk management}

The risk function of Book~16, chapter~12, is led by the chief risk officer and sits outside the businesses it limits.
Its quantitative work is the measurement of Book~6: value at risk and expected shortfall (chapter~21), stress tests
(chapter~22), counterparty exposure (chapter~17), and the capital calculations of chapter~23. The \omterm{def:in:bank-quant:risk}{risk quant} of
chapter~18 builds those models; the risk manager uses them to set limits, watch positions against them, and escalate
when a desk breaches one or when a risk the models do not see appears. The four-eyes principle and the risk committee of
Book~16 are the risk manager's tools of authority.

A risk manager is judged by what does not happen: no unexplained loss, no limit breach unreported, no desk surprised by a
stress. The work needs the quantitative literacy to read a model's output critically and the standing to say no to a
trader who earns more. Professional certification is common and optional: GARP's Financial Risk Manager requires passing
``two multiple-choice Exams'' and ``at least two years of relevant work experience''.

\section{Operations}

Operations makes sure that what the front office traded happens: trades confirmed with counterparties, settled at the
central securities depository, margined at the clearing house, reconciled against the firm's own books and the broker's
(Book~16, chapter~13; Book~15, chapter~22, for the systems). At a trading firm that trades millions of times a day the
work is automated, and the operations analyst's job is the exceptions: a break in a reconciliation, a failed settlement, a
corporate action the systems missed. The quantitative content is data and process: finding the pattern in the breaks,
automating what repeats, measuring the cost of fails.

\section{Product control}

\begin{definition}[Product controller]\label{def:in:risk-operations-product-control-and-compliance:controller}
A \emph{product controller}\index{product controller} is a member of the finance function who produces and explains
each trading desk's daily profit and loss independently of the desk, checks the prices the desk marks its positions at
(independent price verification), and signs off the valuation reserves.
\end{definition}

Product control (Book~6, chapter~27; Book~16, chapter~15) is where a trading loss is found first. Its daily profit and
loss explain attributes the desk's result to market moves, new trades and residuals, and an unexplained residual is the
\omterm{def:in:risk-operations-product-control-and-compliance:controller}{product controller}'s to chase. The quantitative work is the valuation models of Books~5 and~6 read from outside, the
Greeks-based explain, and the independent sources for prices. Product control teaches every desk's products and every
desk's numbers, which makes it a useful first job for a quantitative graduate at a bank.

\section{Compliance}

\begin{definition}[Compliance officer]\label{def:in:risk-operations-product-control-and-compliance:compliance}
A \emph{compliance officer}\index{compliance officer} is a member of the function that makes sure a firm and its staff
follow the laws and rules that apply to them: advises the business, trains and certifies staff, monitors trading and
communications for market abuse, and reports to regulators when required.
\end{definition}

The compliance function of Book~16, chapter~16, has become quantitative through surveillance: the detection of market
abuse (spoofing, layering, insider dealing, marking the close) in order and trade data, which Book~9, chapter~29, and
Book~15, chapter~30, describe as models and systems. In the United Kingdom compliance is also a named responsibility: the
regulator's list of senior management functions includes ``SMF 16 Compliance oversight function'' and ``SMF 17 Money
laundering reporting function'' (chapter~25).

\begin{dated}{2026-09}{Rules that shape the control functions}\label{dat:in:risk-operations-product-control-and-compliance:rules}
EBA Guidelines on sound remuneration policies (EBA/GL/2021/04, 2 July 2021): \omterm{def:in:risk-operations-product-control-and-compliance:control}{control functions}' pay ``should be
predominantly fixed''; their variable pay determined ``separately from the business units they control''; institutions
``should consider setting a significantly lower ratio between the variable and fixed components of remuneration for
\omterm{def:in:risk-operations-product-control-and-compliance:control}{control functions}''. United Kingdom: SMF 16 (compliance oversight) and SMF 17 (money laundering reporting) are
FCA-required senior management functions. GARP's Financial Risk Manager: two exams (100 and 80 questions) and two years of
relevant experience.
\end{dated}

\section{Control functions as careers and as routes to the front office}

Two questions matter to a candidate: what the functions pay against the front office, and whether they lead anywhere.

\begin{dated}{2025-09}{What control functions are paid: the public evidence}\label{dat:in:risk-operations-product-control-and-compliance:pay}
Labour condition applications at banks, fiscal 2025: risk titles 380 applications (5 employers), median offered base
\$136\,776; traders and \omterm{def:in:quant-researcher:def}{quantitative researchers} 883 (9 employers), median \$160\,000; ratio 0.855 (95\% interval
0.807--0.903). At the trading firms pooled: 17 risk filings from 7 employers, median \$193\,003, ratio 0.965
(0.925--1.151). Occupational survey, May 2025, securities industry, median wages: financial risk specialists \$133\,070,
accountants and auditors \$105\,160, \omterm{def:in:risk-operations-product-control-and-compliance:compliance}{compliance officers} \$100\,160, brokerage clerks \$70\,070. EBA \omterm{def:in:pay-levels-by-role-firm-type-and-seniority:median}{high earners} 2024,
credit institutions: 53 in independent \omterm{def:in:risk-operations-product-control-and-compliance:control}{control functions}, variable-to-fixed ratio 0.88, against 1.23 in investment
banking.
\end{dated}

\begin{omfigure}
\begin{tikzpicture}
\begin{axis}[width=10.5cm,height=5.2cm,xmin=0.5,xmax=5.5,ymin=0.5,ymax=1.3,xtick={1,2,3,4,5},
  xticklabels={all levels,I,II,III,IV},xlabel={\small wage level},ylabel={\small control over front office},
  tick label style={font=\footnotesize},grid=major,grid style={gray!20},table/col sep=comma,unbounded coords=jump,
  legend style={legend cell align=left,font=\scriptsize,at={(0.5,-0.3)},anchor=north,legend columns=3,/tikz/every even column/.append style={column sep=8pt}}]
\addplot[only marks,mark=o,black!60,thick,error bars/.cd,y dir=both,y explicit]
  table[x expr=\thisrow{pos}-0.15,y=bank2021,y error plus expr=\thisrow{bank2021_hi}-\thisrow{bank2021},
  y error minus expr=\thisrow{bank2021}-\thisrow{bank2021_lo}]{figdata/industry/24-risk-operations-product-control-and-compliance/ratio.csv};
\addplot[only marks,mark=*,omDef,thick,error bars/.cd,y dir=both,y explicit]
  table[x=pos,y=bank2025,y error plus expr=\thisrow{bank2025_hi}-\thisrow{bank2025},
  y error minus expr=\thisrow{bank2025}-\thisrow{bank2025_lo}]{figdata/industry/24-risk-operations-product-control-and-compliance/ratio.csv};
\addplot[only marks,mark=square*,omThm,thick,error bars/.cd,y dir=both,y explicit]
  table[x expr=\thisrow{pos}+0.15,y=trading2025,y error plus expr=\thisrow{trading2025_hi}-\thisrow{trading2025},
  y error minus expr=\thisrow{trading2025}-\thisrow{trading2025_lo}]{figdata/industry/24-risk-operations-product-control-and-compliance/ratio.csv};
\draw[black!50] (axis cs:0.5,1) -- (axis cs:5.5,1);
\legend{banks 2021,banks 2025,trading firms 2025}
\end{axis}
\end{tikzpicture}

\omcaption{Median offered base of risk-title filings over that of trader and quantitative-researcher filings, with 95\%
bootstrap intervals: banks in fiscal 2021 and 2025, overall and by \omterm{def:in:pay-levels-by-role-firm-type-and-seniority:soc}{wage level}; the trading firms pooled, overall only
(too few filings by level). Data: \texttt{data/industry/lca\_control\_gap.csv}, through
\texttt{in\_control.ratios}.}\label{fig:in:risk-operations-product-control-and-compliance:ratio}
\end{omfigure}

At banks the \omterm{def:in:risk-operations-product-control-and-compliance:control}{control function}'s offered base ran about 15\% below the front office's in fiscal 2025 overall, and 15\% to
32\% below at each level (\cref{fig:in:risk-operations-product-control-and-compliance:ratio}); in fiscal 2021 the two
were level overall. The gap
in base pay is the smaller part of the difference: in the banks' reports of their highest earners, \omterm{def:in:risk-operations-product-control-and-compliance:control}{control functions}'
variable pay is 0.88 of fixed against 1.23 in investment banking, as the guidelines intend
(\cref{fig:in:risk-operations-product-control-and-compliance:eba}). At the trading firms, the few risk filings sit close
to the front office's base, but their interval is wide.

\begin{omfigure}
\begin{tikzpicture}
\begin{axis}[xbar,width=10.5cm,height=5cm,xmin=0,xmax=1.8,ymin=0.4,ymax=7.6,bar width=8pt,ytick=data,
  yticklabels from table={figdata/industry/24-risk-operations-product-control-and-compliance/eba.csv}{area},
  yticklabel style={font=\scriptsize},xticklabel style={font=\footnotesize},
  xlabel={\small variable pay over fixed pay, high earners 2024},grid=major,grid style={gray!20},table/col sep=comma,
  nodes near coords,nodes near coords style={font=\scriptsize,anchor=west,xshift=1pt},point meta=explicit symbolic]
\addplot[fill=omDef,draw=none] table[x=ratio,y=pos,meta=n]
  {figdata/industry/24-risk-operations-product-control-and-compliance/eba.csv};
\end{axis}
\end{tikzpicture}

\omcaption{Ratio of variable to fixed pay of EU credit institutions' staff paid one million euros or more in 2024, by
business area; labels give the number of such staff. Independent \omterm{def:in:risk-operations-product-control-and-compliance:control}{control functions}' ratio is below those of
investment banking, asset management and corporate functions. Data: \texttt{data/industry/eba\_high\_earners.csv} (EBA), through
\texttt{in\_control.eba}.}\label{fig:in:risk-operations-product-control-and-compliance:eba}
\end{omfigure}

The survey shows the same order among the occupations that fill the functions
(\cref{fig:in:risk-operations-product-control-and-compliance:survey}): risk specialists paid near the financial
analysts, accountants and \omterm{def:in:risk-operations-product-control-and-compliance:compliance}{compliance officers} below, operations clerks well below.

\begin{omfigure}
\begin{tikzpicture}
\begin{axis}[width=10.5cm,height=5cm,xmin=20,xmax=260,ymin=0.4,ymax=5.6,ytick=data,
  yticklabels from table={figdata/industry/24-risk-operations-product-control-and-compliance/survey.csv}{label},
  yticklabel style={font=\scriptsize},xticklabel style={font=\footnotesize},xlabel={\small annual wage, \$ thousand},
  grid=major,grid style={gray!20},table/col sep=comma]
\addplot[draw=none,mark=none,error bars/.cd,x dir=both,x explicit,error mark=none,error bar style={omDef,line width=1pt}]
  table[x=p50,y=pos,x error plus expr=\thisrow{p90}-\thisrow{p50},x error minus expr=\thisrow{p50}-\thisrow{p10}]
  {figdata/industry/24-risk-operations-product-control-and-compliance/survey.csv};
\addplot[draw=none,mark=none,error bars/.cd,x dir=both,x explicit,error mark=none,error bar style={omDef!40,line width=7pt}]
  table[x=p50,y=pos,x error plus expr=\thisrow{p75}-\thisrow{p50},x error minus expr=\thisrow{p50}-\thisrow{p25}]
  {figdata/industry/24-risk-operations-product-control-and-compliance/survey.csv};
\addplot[only marks,mark=|,mark size=5pt,line width=1.5pt,omThm] table[x=p50,y=pos]{figdata/industry/24-risk-operations-product-control-and-compliance/survey.csv};
\end{axis}
\end{tikzpicture}

\omcaption{Wages in the securities industry of the occupations that fill the control and support functions, with
financial and investment analysts for comparison, May 2025 occupational survey: 10th to 90th percentile (thin),
interquartile range (thick), median (mark). Data: \texttt{data/industry/oews\_roles.csv}, through
\texttt{in\_control.survey}.}\label{fig:in:risk-operations-product-control-and-compliance:survey}
\end{omfigure}

The routes run both ways. Product control and market risk teach every desk's products and numbers, and people move from
them to the desks they controlled, as strategists, \omterm{def:in:trader:kinds}{risk traders} or \omterm{def:in:structurer-sales-and-sales-trader:structurer}{structurers} (chapters~18 and~23); the reverse move,
from the front office into risk or compliance, takes experience into the function that checks it. Model validation
(chapter~18) sits between. Chapter~28 measures such moves where public data allow. For a candidate the question is less
the first title than what the job teaches and who can see the work: a \omterm{def:in:risk-operations-product-control-and-compliance:control}{control function} sees the whole firm, the front
office sees one desk.

\section{Tutorial: the price of independence}\label{tut:in:risk-operations-product-control-and-compliance:tut}

\begin{tutorial}
\textbf{Goal.} Measure the \omterm{def:in:risk-operations-product-control-and-compliance:control}{control functions}' base pay against the front office's, and their pay mix against the
business areas they control.
\textbf{End state:} \cref{fig:in:risk-operations-product-control-and-compliance:ratio,fig:in:risk-operations-product-control-and-compliance:eba,fig:in:risk-operations-product-control-and-compliance:survey}.
\begin{tutsteps}
\item \textbf{Base pay.} \texttt{in\_lca\_control\_derive.py} compares the risk family with the trader and \omterm{def:in:quant-researcher:def}{quantitative
researcher} families at banks and at the trading firms pooled, with \texttt{firm.roles.median\_gap} (chapter~19), from
chapter~14's cache.
\item \textbf{Pay mix.} \texttt{firm.roles.fixed\_share} and \texttt{mix\_gap} turn the EBA's variable-to-fixed ratios
into shares of fixed pay (\cref{lst:in:risk-operations-product-control-and-compliance:mix}).
\omcode{code/firm/roles/firm_roles.py}{430}{437}{Share of fixed pay and the gap in pay mix between a control function and
a business area.}\label{lst:in:risk-operations-product-control-and-compliance:mix}
\item \textbf{Survey.} \texttt{oews\_roles.csv} for the functions' occupations in the securities industry.
\end{tutsteps}
\end{tutorial}

For the \omterm{def:in:pay-levels-by-role-firm-type-and-seniority:median}{high earners}: fixed pay is 53.3\% of the \omterm{def:in:risk-operations-product-control-and-compliance:control}{control functions}' total and 44.8\% of investment banking's; the ratio of
their variable-to-fixed ratios is 0.71.

\textbf{What to change next.} Separate compliance and product control titles in the filings, which needs the job titles
this book's single pass did not keep; add bonus from the banks' remuneration disclosures by business line; follow
people across roles in public profiles (chapter~28).

\section{Build: control cards and pay mix}\label{bld:in:risk-operations-product-control-and-compliance:roles}

\begin{build}
\bfield{Purpose} Add the control and support functions to the registry and measure the pay mix that the rules require.
\bfield{Interface} \texttt{firm.roles}: the cards \texttt{risk manager}, \texttt{operations analyst}, \texttt{product
controller}, \texttt{compliance officer}; \texttt{fixed\_share(v)}; \texttt{mix\_gap(control\_ratio, business\_ratio)};
with \texttt{median\_gap} for base pay.
\bfield{Rules} A ratio of variable to fixed pay is converted, never averaged, across areas; the base-pay comparison keeps
the three-employer rule.
\bfield{Acceptance tests} \texttt{code/firm/roles/tests/}: equal variable and fixed pay give a half; no variable pay gives
one; a control ratio half the business's gives a ratio of ratios of one half.
\bfield{Stretch} A bootstrap over employers rather than filings; the effect of a \omterm{def:in:pay-regulation-and-tax-by-location:cap}{bonus cap} (chapter~15) on the mix.
\end{build}

\begin{omsources}
\item European Banking Authority (2021), Guidelines on sound remuneration policies under Directive 2013/36/EU,
EBA/GL/2021/04; EBA, high earners 2024.
\item FCA Handbook, SUP 10C; GARP, Financial Risk Manager.
\item Chapter~14's tables from the Department of Labor's LCA files; BLS occupational survey, May 2025.
\item Book~16, chapters~12--16; Book~6, chapters~21--28.
\end{omsources}

\section{Exercises}

\begin{exercise}[$\star$]\label{exo:in:risk-operations-product-control-and-compliance:1}
By how much, in dollars, did the banks' median offer for risk titles fall short of the front office's in fiscal 2025?
\end{exercise}

\begin{exercise}[$\star$]\label{exo:in:risk-operations-product-control-and-compliance:2}
A control-function \omterm{def:in:pay-levels-by-role-firm-type-and-seniority:median}{high earner}'s variable pay is 0.88 of fixed. What share of total pay is fixed?
\end{exercise}

\begin{exercise}[$\star$]\label{exo:in:risk-operations-product-control-and-compliance:3}
Which UK senior management functions name the compliance and anti-money-laundering responsibilities?
\end{exercise}

\begin{exercise}[$\star\star$]\label{exo:in:risk-operations-product-control-and-compliance:4}
The banks' control-to-front ratio was 1.01 in fiscal 2021 and 0.86 in fiscal 2025. Using the medians (\$131\,720 and
\$130\,000 in 2021; \$136\,776 and \$160\,000 in 2025), which side moved?
\end{exercise}

\begin{exercise}[$\star\star$]\label{exo:in:risk-operations-product-control-and-compliance:5}
Why would a regulator want \omterm{def:in:risk-operations-product-control-and-compliance:control}{control functions}' pay to be predominantly fixed, and what does the rule cost?
\end{exercise}

\begin{exercise}[$\star\star$]\label{exo:in:risk-operations-product-control-and-compliance:6}
A graduate chooses between product control and a desk-analyst job at the same bank. What does each teach, and which is
easier to leave?
\end{exercise}

\begin{exercise}[$\star\star\star$]\label{exo:in:risk-operations-product-control-and-compliance:7}
\emph{Coding.} The trading firms' ratio rests on 17 filings from 7 employers. Explain why its interval is asymmetric
(0.925 to 1.151 around 0.965), and what a bootstrap over employers rather than filings would do to it.
\end{exercise}

\begin{exercise}[$\star\star\star$]\label{exo:in:risk-operations-product-control-and-compliance:8}
\emph{Find the flaw.} ``Risk managers at banks are paid 15\% less than traders, so the banks do not value risk
management.''
\end{exercise}

\section{Problem: The Price of Independence}

\begin{problem}[{Weekend problem --- the price of independence}]\label{pb:in:risk-operations-product-control-and-compliance:1}
A quantitative graduate has offers in market risk at a bank and as a junior researcher at a market maker, and asks what
choosing the \omterm{def:in:risk-operations-product-control-and-compliance:control}{control function} costs.

\medskip\noindent\textbf{Part I --- The functions.}
\begin{enumerate}
\item Define a \omterm{def:in:risk-operations-product-control-and-compliance:control}{control function}, a \omterm{def:in:risk-operations-product-control-and-compliance:controller}{product controller} and a \omterm{def:in:risk-operations-product-control-and-compliance:compliance}{compliance officer}.
\item Which functions does the EBA list as independent \omterm{def:in:risk-operations-product-control-and-compliance:control}{control functions}?
\item What does a risk manager do with the \omterm{def:in:bank-quant:risk}{risk quant}'s models?
\item What does operations do at a firm that trades millions of times a day?
\item How has compliance become quantitative?
\end{enumerate}

\medskip\noindent\textbf{Part II --- The rules.}
\begin{enumerate}[resume]
\item State the EBA guidelines' rules on \omterm{def:in:risk-operations-product-control-and-compliance:control}{control functions}' pay.
\item Why must variable pay be set separately from the businesses controlled?
\item What are SMF 16 and SMF 17?
\item What does the FRM require?
\item What do the rules imply for a control-function career's pay path?
\end{enumerate}

\medskip\noindent\textbf{Part III --- The evidence.}
\begin{enumerate}[resume]
\item Give the banks' control-to-front ratio overall and by level in fiscal 2025.
\item Give the trading firms' ratio and say why it is uncertain.
\item Compare the \omterm{def:in:pay-levels-by-role-firm-type-and-seniority:median}{high earners}' pay mix in \omterm{def:in:risk-operations-product-control-and-compliance:control}{control functions} and investment banking.
\item Rank the functions' occupations by the survey's medians.
\item How did the ratio change since fiscal 2021?
\end{enumerate}

\medskip\noindent\textbf{Part IV --- The verdict.}
\begin{enumerate}[resume]
\item State the \emph{named result}: the ratio of control-function to front-office median offered base at banks and at
trading firms, with bootstrap intervals.
\item What does base pay miss here?
\item What does the market-risk job teach that the research job does not?
\item What routes lead from market risk to the front office?
\item In two sentences, answer the graduate.
\end{enumerate}
\end{problem}

\section{Interview questions}

\begin{interviewq}[$\star$ \iqroles{risk}]\label{iq:in:risk-operations-product-control-and-compliance:1}
A desk's P\&L explain shows a large unexplained residual today. What do you check first?
\end{interviewq}

\begin{interviewq}[$\star$ \iqroles{risk}]\label{iq:in:risk-operations-product-control-and-compliance:2}
A trader breaches a limit by 5\% for ten minutes and says the model is wrong. What do you do?
\end{interviewq}

\begin{interviewq}[$\star\star$ \iqroles{risk, bank}]\label{iq:in:risk-operations-product-control-and-compliance:3}
How would you verify the price of a position with no liquid market quote?
\end{interviewq}

\begin{interviewq}[$\star\star$ \iqroles{risk, developer}]\label{iq:in:risk-operations-product-control-and-compliance:4}
Design a detector for spoofing in order data. What is its false-positive problem?
\end{interviewq}

\begin{interviewq}[$\star\star$ \iqroles{risk}]\label{iq:in:risk-operations-product-control-and-compliance:5}
Your value at risk fell sharply while the positions did not change. Why might that be?
\end{interviewq}

\begin{interviewq}[$\star\star\star$ \iqroles{risk, trader}]\label{iq:in:risk-operations-product-control-and-compliance:6}
Two desks each hedge the other's risk and both show profits. What would make you suspicious, and how would you check?
\end{interviewq}
